\documentclass[9.5pt,a4paper]{extarticle}

% --- Packages ---
\usepackage[utf8]{inputenc}
\usepackage[T1]{fontenc}
\usepackage[spanish,es-tabla]{babel}
\usepackage[a4paper,margin=0.42in,top=0.38in,bottom=0.38in]{geometry}
\usepackage{multicol}
\usepackage{enumitem}
\usepackage{xcolor}
\usepackage{titlesec}
\usepackage{booktabs}
\usepackage{tcolorbox}
\usepackage{hyperref}
\usepackage{amsmath,amssymb}
\usepackage{microtype}

% --- Black & White / Grayscale Palette ---
\definecolor{boxbg}{gray}{0.96}
\definecolor{boxframe}{gray}{0.15}

\hypersetup{
    colorlinks=true,
    linkcolor=black,
    citecolor=black,
    urlcolor=black,
    pdftitle={IA Generativa - Entrega 1: Asignacion de Turnos Medicos},
    pdfauthor={Grupo 11}
}

% --- Section Spacing & Formatting ---
\titleformat{\section}{\fontsize{9.8}{11.2}\bfseries\scshape}{}{0em}{}[\vspace{-3pt}\rule{\linewidth}{0.65pt}\vspace{-2.5pt}]
\titlespacing*{\section}{0pt}{4pt}{2.5pt}

\titleformat{\subsection}{\fontsize{8.8}{10}\bfseries}{}{0em}{}
\titlespacing*{\subsection}{0pt}{2.5pt}{1.5pt}

\setlength{\columnsep}{14pt}
\setlength{\parskip}{1.8pt}
\setlength{\parindent}{0pt}
\setlist{noitemsep,topsep=1.5pt,partopsep=0pt,parsep=0.5pt,leftmargin=1.1em}

\pagestyle{empty} % Documento de 1 sola pagina

\begin{document}

% --- Header Banner (Black & White) ---
\begin{tcolorbox}[colback=boxbg,colframe=boxframe,arc=1.5mm,boxrule=0.8pt,top=3pt,bottom=3pt,left=8pt,right=8pt]
\begin{minipage}[c]{0.68\textwidth}
    {\fontsize{12}{13.5}\selectfont \textbf{Asignación de Turnos Médicos con LLMs Open-Weight}}\\[1.5pt]
    {\fontsize{8.5}{9.8}\selectfont \textbf{Deliverable 1 \textbar\ Generative Artificial Intelligence (580694)}}
\end{minipage}
\hfill
\begin{minipage}[c]{0.30\textwidth}
    \raggedleft
    {\fontsize{8}{9}\selectfont \textbf{Universidad de Concepción}}\\
    {\fontsize{8.5}{9.5}\selectfont \textbf{Grupo 11}}\\
    {\fontsize{7}{8}\selectfont \textbf{Repo:} \href{https://github.com/matirocha/genai-grupo11}{\texttt{github.com/matirocha/genai-grupo11}}}
\end{minipage}
\end{tcolorbox}

\vspace{-2pt}

\begin{multicols}{2}

% ==================== SECCIÓN 1 ====================
\section{1. Definición de la Tarea y Ground Truth}
\textbf{Contexto:} El problema de asignación de turnos médicos (\textit{Medical Staff Scheduling Problem}, MSSP) es un desafío combinatorio NP-hard de satisfacción de restricciones (CSP) crítico en la gestión hospitalaria.

\textbf{Definición de la Tarea:} Dado un conjunto de personal médico $S = \{s_1, \dots, s_N\}$ (con especialidades, horas máximas semanales por contrato $H_i$ e indisponibilidades solicitadas $U_i$) y la demanda semanal $D = \{d_1, \dots, d_{21}\}$ en 3 turnos de 8h (\textit{Mañana, Tarde, Noche}), el modelo debe generar la matriz de asignación completa en formato JSON válido.

\textbf{Ground Truth y Evaluador Programático:}
La salida se evalúa mediante un oráculo determinista en Python (\texttt{src/verifier.py}) que verifica 0 violaciones en 6 Restricciones Duras (\textit{Hard Constraints}, HC):
\begin{itemize}[label=$\blacktriangleright$]
    \item \textbf{HC1 (Descanso Obligatorio):} $\ge 16$h continuas de descanso entre turnos (un turno Noche no puede ser seguido por Mañana o Tarde al día siguiente).
    \item \textbf{HC2 (Turno Único Diario):} Máximo 1 turno (8h) por médico al día.
    \item \textbf{HC3 (Demanda de Dotación):} Cantidad exacta de personal por turno.
    \item \textbf{HC4 (Cobertura de Especialidad):} Cuota mínima de especialistas (ej. $\ge 1$ Anestesiólogo en Noche, $\ge 1$ Urgenciólogo en Día).
    \item \textbf{HC5 (Límite Legal Semanal):} Horas totales asignadas $\le H_i$ ($\le 40$h).
    \item \textbf{HC6 (Indisponibilidad):} Cero asignaciones en franjas bloqueadas.
\end{itemize}

% ==================== SECCIÓN 2 ====================
\section{2. Diagnóstico del Fallo y Evidencia}
El \textit{prompting} directo (zero-shot) en modelos pequeños ($\le 8\text{B}$) falla sistemáticamente debido a tres limitaciones estructurales:

\begin{enumerate}[label=\textbf{\arabic*.}]
    \item \textbf{Decisión Voraz vs. Búsqueda Global CSP:}
    Los LLMs generan el horario token a token de lunes a domingo. Al carecer de mecanismos de retroceso (\textit{backtracking}), las decisiones iniciales agotan vorazmente las horas de especialistas, forzando violaciones inevitables hacia el final de la semana (viernes a domingo).
    \item \textbf{Deriva en el Seguimiento de Estado y Aritmética:}
    El modelo no mantiene contadores acumulados exactos de horas trabajadas ($\sum h_i$), generando sobrecargas ilegales ($\text{HC5}$) y duplicidades ($\text{HC2}$).
    \item \textbf{Alucinación de Esquema e Identificadores:}
    En salidas largas, los modelos pequeños suelen truncar claves JSON, repetir IDs en un mismo turno o inventar personal inexistente.
\end{enumerate}

\textbf{Evidencia Empírica de Línea Base:} En pruebas zero-shot sobre 10 médicos y 21 turnos, los modelos de 3B obtuvieron un \textbf{0\% de horarios válidos}, con un promedio de \textbf{7.4 violaciones duras por intento} (principalmente descanso HC1 y exceso de horas HC5).

% ==================== SECCIÓN 3 ====================
\section{3. Modelos Candidatos y Justificación}
Para abordar la tarea con alta eficiencia computacional y rápida iteración, seleccionamos tres modelos compactos de última generación en el rango de \textbf{3.0B a 3.8B} (dentro del límite de 8B):

\begin{center}
\vspace{-2pt}
\resizebox{\columnwidth}{!}{%
\begin{tabular}{@{}lcccp{3.8cm}@{}}
\toprule
\textbf{Modelo} & \textbf{Params} & \textbf{Ctx} & \textbf{GSM8k} & \textbf{Justificación Arquitectural} \\
\midrule
\textbf{Qwen 2.5 3B-Instruct} & 3.09B & 128k & 84.5\% & SOTA en JSON estructurado y seguimiento de restricciones lógicas. \\
\textbf{Phi-3.5-mini-Instruct} & 3.82B & 128k & 86.0\% & Datos sintéticos de alto razonamiento; sobresale en lógica paso a paso. \\
\textbf{Llama 3.2 3B-Instruct} & 3.21B & 128k & 77.2\% & Modelo ligero y veloz; óptimo para fine-tuning eficiente con LoRA. \\
\bottomrule
\end{tabular}%
}
\vspace{-2pt}
\end{center}

\textit{Justificación de Escala Compacta:} El rango de $\approx 3\text{B}$ ofrece latencia sub-segundo, mínimo consumo de memoria para fine-tuning (PEFT/LoRA) y permite demostrar que marcos de razonamiento estructurado (CoT / Tool-calling / Verificador) resuelven eficazmente problemas CSP complejos.

% ==================== SECCIÓN 4 ====================
\section{4. Factibilidad de Ejecución}
La infraestructura experimental está 100\% garantizada tanto en entorno local como en la nube:
\begin{itemize}[label=$\blacktriangleright$]
    \item \textbf{Hardware Local Dedicado:} Estación de trabajo equipada con \textbf{AMD Ryzen 7 7700}, GPU \textbf{Nvidia RTX 5060 (8 GB VRAM)} y \textbf{32 GB RAM DDR5}. Soporta inferencia nativa en FP16 ($\approx 6.5\text{ GB}$) y cuantizada en 4-bit NF4 ($\approx 2.4\text{ GB}$), además de fine-tuning LoRA rápido sin límites de sesión.
    \item \textbf{Google Colab Gratuito (GPU Nvidia T4, 15 GB VRAM):} Entorno en la nube 100\% reproducible; ocupa $\approx 2.4\text{ GB}$ en 4-bit NF4, dejando $>12\text{ GB}$ libres para memoria de contexto KV.
    \item \textbf{Rendimiento y Latencia:} Alcanza $\approx 35\text{--}45\text{ tokens/s}$, completando la generación del horario semanal en menos de 20 segundos.
\end{itemize}

% ==================== SECCIÓN 5 ====================
\section{5. Plan Semestral y Repositorio}
\textbf{Estrategia de Solución Semestral:} Atacaremos el fallo mediante:
\textbf{(1)} \textit{Chain-of-Thought} multietapa con verificación intermedia, \textbf{(2)} \textit{Tool-Calling} integrando oráculos y \textit{solvers} CSP (Z3/OR-Tools), y \textbf{(3)} \textit{Fine-Tuning} LoRA supervisado sobre trazas válidas de resolución.

\textbf{Estructura del Repositorio:} Código reproducible en GitHub:
\begin{itemize}[label=$\blacktriangleright$]
    \item \texttt{/data}: Esquemas JSON de personal y demanda de turnos.
    \item \texttt{/src}: \texttt{verifier.py} (\textit{Ground Truth}) y \texttt{baseline\_direct.py}.
    \item \texttt{/notebooks}: \texttt{Colab\_Feasibility.ipynb} con perfilado de VRAM.
\end{itemize}

\end{multicols}

\end{document}
